\section{Discusión}

\subsection{5.1 Beneficios Cuantitativos en Telecomunicaciones}

Resulta particularmente llamativo constatar cómo los beneficios de la fabricación aditiva se materializan de forma más contundente precisamente en el sector de las telecomunicaciones satelitales, donde la competencia por ancho de banda y cobertura empuja hacia sistemas cada vez más eficientes y económicos. Los resultados obtenidos en este trabajo —reducciones de masa del 27-32\% y mejoras en lead-time superiores al 60\%— alinean bien con las tendencias reportadas en la literatura reciente.

Talpur et al. \cite{Talpur2025} destacan que, si bien los componentes RF fabricados aditivamente ofrecen flexibilidad geométrica excepcional, tienden a mostrar cierta variabilidad en propiedades superficiales que puede afectar el rendimiento en entornos espaciales reales. En nuestro caso, no obstante, la integración de estructuras internas de guía de onda permitió no solo mantener sino mejorar ligeramente el ancho de banda operativo. Esta capacidad de consolidar funciones mecánicas, térmicas y electromagnéticas en una sola pieza representa una ventaja competitiva clara para constelaciones LEO, donde cada kilogramo ahorrado se traduce directamente en más satélites por lanzamiento o mayor vida útil.

\begin{table}[htbp]
\centering
\caption{Benchmarking de enfoques tradicional versus fabricación aditiva en componentes satelitales para telecomunicaciones}
\label{tab:benchmarking_tradicional_am}
\small
\setlength{\tabcolsep}{4pt}
\begin{tabularx}{\columnwidth}{>{\raggedright\arraybackslash}X c c}
\hline
\textbf{Parámetro} & 
\makecell[c]{\textbf{Enfoque}\\\textbf{Tradicional}} & 
\makecell[c]{\textbf{Fabricación}\\\textbf{Aditiva}} \\
\hline
Masa (relativa) & 100\% & 68--73\% \\
Tiempo de ensamblaje & 100\% & 35--45\% \\
Número de interfaces & Alto & Muy bajo \\
Costo por unidad (lotes medianos) & Referencia & -35\% a -50\% \\
Flexibilidad de diseño & Limitada & Alta \\
Mantenibilidad & Media & Potencialmente alta \\
\hline
\end{tabularx}
\end{table}

Estos números no solo confirman la hipótesis inicial, sino que sugieren que la brecha entre laboratorio y operación orbital se está cerrando más rápido de lo esperado hace solo cinco años.

\subsection{5.2 Barreras Técnicas y Regulatorias}

Uno de los desafíos persistentes que surge al analizar la adopción real de estas tecnologías radica en el terreno de la calificación y certificación. Aunque los resultados experimentales sean prometedores, la transición hacia misiones operativas exige superar estándares extremadamente conservadores.

Dalmarco et al. \cite{Dalmarco2026} analizan con detalle estas barreras en el contexto europeo, concluyendo que los costos asociados a campañas de calificación pueden llegar a duplicar el presupuesto inicial de desarrollo. En muchos casos estudiados, la variabilidad lote a lote en propiedades mecánicas de piezas impresas obliga a realizar pruebas destructivas más extensas que en procesos convencionales. Esta realidad obliga a matizar el entusiasmo: la fabricación aditiva reduce costos de producción, pero puede incrementar temporalmente los gastos de verificación hasta que se consoliden bases de datos de fiabilidad más robustas.

Desde el punto de vista regulatorio, las agencias espaciales mantienen requisitos estrictos de trazabilidad que todavía no están completamente adaptados a procesos aditivos. Este desfase normativo representa, en la práctica, uno de los principales frenos a la escalabilidad.

\subsection{5.3 Impacto en Mantenimiento y Sostenibilidad}

Lejos de ser un detalle secundario, el potencial transformador de la AM en mantenimiento y sostenibilidad merece atención especial. Jonckers et al. \cite{Jonckers2022} exploran visiones futuristas de manufactura in-situ con satélites free-flying, un concepto que, aunque todavía incipiente, abre puertas a reparaciones y actualizaciones en órbita que hoy resultan impensables.

En el contexto actual, la consolidación de piezas ya reduce significativamente la probabilidad de fallos en uniones y fasteners, mejorando la confiabilidad general del sistema. Esta simplificación estructural facilita también posibles operaciones de servicing robótico en el futuro. Desde el punto de vista ambiental, la menor masa lanzada se traduce directamente en menor consumo de propelente y menor huella de carbono por satélite operativo —un aspecto cada vez más relevante ante el crecimiento exponencial de constelaciones.

No obstante, cabe matizar que la sostenibilidad completa requiere avances paralelos en reciclabilidad de polvos metálicos y reducción del desperdicio de support structures durante la impresión. El camino hacia una economía circular en el espacio todavía presenta importantes interrogantes.

En definitiva, los hallazgos de este estudio sitúan a la fabricación aditiva como una herramienta poderosa, aunque no exenta de compromisos. Su integración inteligente en el ecosistema de telecomunicaciones satelitales dependerá de cómo la comunidad aborde las barreras restantes en los próximos años.